%-----------------------------------------------------------------------------------------------
\makeatletter
\immediate\write18{datelog > \jobname.info} % site script for $(date '+%Y-%m-%d %Hh%Mm%Ss')
\makeatother
%-----------------------------------------------------------------------------------------------
%-----------------------------------------------------------------------------------------------
\usetheme{Copenhagen}
\usepackage{beamercolorthemeCNThermSci2}
\usefonttheme{serif}
%-----------------------------------------------------------------------------------------------

%-----------------------------------------------------------------------------------------------
%-----------------------------------------------------------------------------------------------
\usetheme{Copenhagen}
\usepackage{beamercolorthemeUTF2}
\usefonttheme{serif}
%-----------------------------------------------------------------------------------------------
\usepackage[utf8]{inputenc}
\usepackage[greek,french,english,brazil]{babel} % last becomes the active one
\usepackage{pslatex}
\usepackage{amssymb,amsmath}
\usepackage{soul}
\usepackage[squaren,Gray,cdot]{SIunits}
\usepackage[nice]{nicefrac}
\usepackage{tikz}
\usepackage{amscd}
\usepackage{stmaryrd}
\usepackage{scalerel}
\usepackage{xspace}
%-----------------------------------------------------------------------------------------------


%-----------------------------------------------------------------------------------------------
%-----------------------------------------------------------------------------------------------
% Mathematical
%-----------------------------------------------------------------------------------------------
\newcommand{\vet}[1]{\underline{{#1}}}
\newcommand{\mat}[1]{\underline{\underline{{#1}}}}
\newcommand{\cub}[1]{\underline{\underline{\underline{{#1}}}}}
\newcommand{\eqdef}{{\ensuremath\stackrel{\text{\tiny def}}{=}}}
%-----------------------------------------------------------------------------------------------
% Linguistic
%-----------------------------------------------------------------------------------------------
\newcommand{\GRtxt}[1]{\begin{otherlanguage}{greek}{{#1}}\end{otherlanguage}}
\newcommand{\FRtxt}[1]{\begin{otherlanguage}{french}{{#1}}\end{otherlanguage}}
%-----------------------------------------------------------------------------------------------
% Presentation
%-----------------------------------------------------------------------------------------------
\newcommand{\BkgImgH}[1]{% Places an image centered on the slide background filling the height
    \usebackgroundtemplate{\parbox{\paperwidth}{%
        \vspace*{1sp}\centering\includegraphics[height=\paperheight]{{#1}}
}}}
\newcommand{\BkgImgW}[1]{% Places an image centered on the slide background filling the width
    \usebackgroundtemplate{\parbox{\paperwidth}{%
        \vspace*{1sp}\centering\includegraphics[width=\paperwidth]{{#1}}
}}}
\newcommand{\ArtEndH}[3]{% Transitions to plain image (last) slide: #1:prefix #2,#3:extensions
    \BkgImgH{root/../art/#1.#2}
    \frame<handout:0>[plain]{%
        \transdissolve\vspace*{72mm}\color{white}\scriptsize\bf\input{root/../art/#1.#3}}
    \usebackgroundtemplate{\mbox{~}}
}
\newcommand{\ArtEndW}[3]{% Transitions to plain image (last) slide: #1:prefix #2,#3:extensions
    \BkgImgW{root/../art/#1.#2}
    \frame<handout:0>[plain]{%
        \transdissolve\vspace*{72mm}\color{white}\scriptsize\bf\input{root/../art/#1.#3}}
    \usebackgroundtemplate{\mbox{~}}
}
\newcommand{\ImgColW}[3]{% Inserts a full-width image in a column
    \includegraphics[width=\columnwidth]{root/../art/#1.#2}\\[-0.5\baselineskip]
    \parbox{\columnwidth}{\tiny\hfill\scalebox{0.85}{\input{root/../art/#1.#3}}}
}
\newcommand{\txtpic}[1]{%
    \fcolorbox{lightgray}{white!90!black}{{#1}} 
}
%-----------------------------------------------------------------------------------------------


%-----------------------------------------------------------------------------------------------
\title{A.07.01 -- Relações de Propriedades Termodinâmicas}
\subtitle{Relações Gerais}
\author{Prof.~C.~Naaktgeboren, PhD}
\date{{\scriptsize\tt%
    \includegraphics[height=6.0mm]{cc/by-nc-nd-88x31.pdf}\\[\smallskipamount]
    https://github.com/CNThermSci/ApplThermSci\\
    Compiled on \input{\jobname.info}
}}
%-----------------------------------------------------------------------------------------------
\begin{document}
%-----------------------------------------------------------------------------------------------
\logo{%
    \parbox{158mm}{% There's a 1mm gap on each side of the 160mm x 90mm slide logo line
        \mode<beamer>{
            \includegraphics[height=6.0mm]{root/00-res/UTFPR/UTFPR-logo-D.pdf}\hfill%
            \includegraphics[height=9.0mm]{root/00-res/logo/CNThermSci-logo-A.pdf}%
        }
        \mode<handout>{
            \includegraphics[height=6.0mm]{root/00-res/UTFPR/UTFPR-logo-W.pdf}\hfill%
            \includegraphics[height=9.0mm]{root/00-res/logo/CNThermSci-logo-W.pdf}%
        }
    }
} % The (delineated, alpha), or washed-out logos
%-----------------------------------------------------------------------------------------------
\frame{\titlepage}
%-----------------------------------------------------------------------------------------------

%-----------------------------------------------------------------------------------------------
\frame{\tableofcontents}
%-----------------------------------------------------------------------------------------------

%-----------------------------------------------------------------------------------------------
\section{Relações das Propriedades $u$, $h$, $s$, $c_P$ e $c_v$}
%-----------------------------------------------------------------------------------------------

%-----------------------------------------------------------------------------------------------
\subsection{Relações para $du$, $dh$ e $ds$}
%-----------------------------------------------------------------------------------------------

    % !j 96 -i8
    %-------------------------------------------------------------------------------------------
    \begin{frame}[allowframebreaks]{Relações para a Energia Interna --}\vspace*{-0em}
        Tomando-se \alert{$u\!:\!u(T, v)$}, tem-se termos em \alert{$dT$} e em \alert{$dv$}:

        \begin{align*}
            du  &= 
                \alert{\left(\frac{\partial u}{\partial T}\right)_v\,dT} +
                \left(\frac{\partial u}{\partial v}\right)_T\,dv
            \qquad\rightharpoondown
            \\[\medskipamount]
            du  &=
                \alert{c_v\,dT} +
                \left(\frac{\partial u}{\partial v}\right)_T\,dv.
        \end{align*}

        Tomando-se \alert{$s\!:\!s(T, v)$}, tem-se termos em \alert{$dT$} e em \alert{$dv$}:

        \begin{align*}
            ds  &= 
                \left(\frac{\partial s}{\partial T}\right)_v\,dT +
                \left(\frac{\partial s}{\partial v}\right)_T\,dv.
        \end{align*}

        \pagebreak
        Tomando \alert{$du = Tds - Pdv$} e agrupando termos em \alert{$dT$} e em \alert{$dv$}:

        \begin{align*}
            du  &=
                T\left(\frac{\partial s}{\partial T}\right)_v\,\alert{dT} +
                \left[
                    T\left(\frac{\partial s}{\partial v}\right)_T - P
                \right]\,\alert{dv},
            \qquad\mbox{que é}
            \\[\medskipamount]
            du  &=
                c_v\,\alert{dT} +
                \left(\frac{\partial u}{\partial v}\right)_T\,\alert{dv}.
        \end{align*}

        \pagebreak
        Igualando os coeficientes de \alert{$dT$} e \alert{$dv$}:

        \begin{align*}
            \alert{\left(\frac{\partial s}{\partial T}\right)_v}
            &\alert{= \frac{c_v}{T}},
            \\[\medskipamount]
            \left(\frac{\partial u}{\partial v}\right)_T
            &= T\left(\frac{\partial s}{\partial v}\right)_T - P.
        \end{align*}

        \pagebreak
        Finalmente, pela Equação de Maxwell baseada na energia de Helmholtz, tem-se:

        \begin{align*}
            \left(\frac{\partial u}{\partial v}\right)_T
            &= T\alert{\left(\frac{\partial s}{\partial v}\right)_T} - P,
            \\[\medskipamount]
            &= T\alert{\left(\frac{\partial P}{\partial T}\right)_v} - P.
        \end{align*}

        O que conduz a

        \begin{equation*}
            \alert{
                du = c_v\,dT +
                \left[
                    T\left(
                        \frac{\partial P}{\partial T}
                    \right)_v - P
                \right]\,dv
            }.
        \end{equation*}

        \pagebreak
        Assim, a variação de \alert{$\Delta u_{12} = u_2 - u_1$} entre estados \alert{$(T_2,
        v_2)$} e \alert{$(T_1, v_1)$} é: \vspace*\bigskipamount

        \begin{equation*}
            \alert{
                \Delta u_{12} = u_2 - u_1 =
                    \int_{T_1}^{T_2}
                        c_v\,dT +
                    \int_{v_1}^{v_2}
                        \left[
                            T\left(
                                \frac{\partial P}{\partial T}
                            \right)_v - P
                        \right]\,dv
            }.
        \end{equation*}

    \end{frame}
    %-------------------------------------------------------------------------------------------

    % !j 96 -i8
    %-------------------------------------------------------------------------------------------
    \begin{frame}[allowframebreaks]{Relações para a Entalpia --}\vspace*{-0em}
        Tomando-se \alert{$h\!:\!h(T, P)$}, tem-se termos em \alert{$dT$} e em \alert{$dP$}:

        \begin{align*}
            dh  &= 
                \alert{\left(\frac{\partial h}{\partial T}\right)_P\,dT} +
                \left(\frac{\partial h}{\partial P}\right)_T\,dP
            \qquad\rightharpoondown
            \\[\medskipamount]
            dh  &=
                \alert{c_P\,dT} +
                \left(\frac{\partial h}{\partial P}\right)_T\,dP
        \end{align*}

        Tomando-se \alert{$s\!:\!s(T, P)$}, tem-se termos em \alert{$dT$} e em \alert{$dP$}:

        \begin{align*}
            ds  &= 
                \left(\frac{\partial s}{\partial T}\right)_P\,dT +
                \left(\frac{\partial s}{\partial P}\right)_T\,dP.
        \end{align*}

        \pagebreak
        Tomando \alert{$dh = Tds + vdP$} e agrupando termos em \alert{$dT$} e em \alert{$dP$}:

        \begin{align*}
            dh  &=
                T\left(\frac{\partial s}{\partial T}\right)_P\,\alert{dT} +
                \left[
                    T\left(\frac{\partial s}{\partial P}\right)_T + v
                \right]\,\alert{dP},
            \qquad\mbox{que é}
            \\[\medskipamount]
            dh  &=
                c_P\,\alert{dT} +
                \left(\frac{\partial h}{\partial P}\right)_T\,\alert{dP}.
        \end{align*}

        \pagebreak
        Igualando os coeficientes de \alert{$dT$} e \alert{$dP$}:

        \begin{align*}
            \alert{\left(\frac{\partial s}{\partial T}\right)_P}
            &\alert{= \frac{c_P}{T}},
            \\[\medskipamount]
            \left(\frac{\partial h}{\partial P}\right)_T
            &= T\left(\frac{\partial s}{\partial P}\right)_T + v.
        \end{align*}

        \pagebreak
        Finalmente, pela Equação de Maxwell baseada na energia de Gibbs, tem-se:

        \begin{align*}
            \left(\frac{\partial h}{\partial P}\right)_T
            &= T\alert{\left(\frac{\partial s}{\partial P}\right)_T} + v,
            \\[\medskipamount]
            &= v \alert{-} T\alert{\left(\frac{\partial v}{\partial P}\right)_T}.
        \end{align*}

        O que conduz a

        \begin{equation*}
            \alert{
                dh = c_P\,dT +
                \left[
                    v -
                    T\left(
                        \frac{\partial v}{\partial T}
                    \right)_P
                \right]\,dP
            }.
        \end{equation*}

        \pagebreak
        Assim, a variação de \alert{$\Delta h_{12} = h_2 - h_1$} entre estados \alert{$(T_2,
        P_2)$} e \alert{$(T_1, P_1)$} é: \vspace*\bigskipamount

        \begin{equation*}
            \alert{
                \Delta h_{12} = h_2 - h_1 =
                    \int_{T_1}^{T_2}
                        c_p\,dT +
                    \int_{P_1}^{P_2}
                        \left[
                            v - T\left(
                                \frac{\partial v}{\partial T}
                            \right)_P
                        \right]\,dP
            }.
        \end{equation*}

    \end{frame}
    %-------------------------------------------------------------------------------------------

    % !j 96 -i8
    %-------------------------------------------------------------------------------------------
    \begin{frame}[allowframebreaks]{Relações para a Entropia --}\vspace*{-0em}
        Tomando-se \alert{$s\!:\!s(T, v)$}, como na dedução de \alert{$du$} e a relação de
        Maxwell com base na energia de \alert{Helmholtz}, tem-se: 

        \begin{equation*}
            \alert{
                ds = 
                    \frac{c_v}{T}\,dT +
                    \left(
                        \frac{\partial P}{\partial T}
                    \right)_v\,dv
            },
        \end{equation*}

        \noindent e assim:

        \begin{equation*}
            \alert{
                \Delta s_{12} = s_2 - s_1 =
                    \int_{T_1}^{T_2}
                        \frac{c_v}{T}\,dT +
                    \int_{v_1}^{v_2}
                        \left(
                            \frac{\partial P}{\partial T}
                        \right)_v\,dv
            }.
        \end{equation*}

        \pagebreak
        Tomando-se \alert{$s\!:\!s(T, P)$}, como na dedução de \alert{$dh$} e a relação de
        Maxwell com base na energia de \alert{Gibbs}, tem-se: 

        \begin{equation*}
            \alert{
                ds = 
                    \frac{c_p}{T}\,dT -
                    \left(
                        \frac{\partial v}{\partial T}
                    \right)_P\,dP
            },
        \end{equation*}

        \noindent e assim:

        \begin{equation*}
            \alert{
                \Delta s_{12} = s_2 - s_1 =
                    \int_{T_1}^{T_2}
                        \frac{c_p}{T}\,dT -
                    \int_{P_1}^{P_2}
                        \left(
                            \frac{\partial v}{\partial T}
                        \right)_P\,dP
            }.
        \end{equation*}

    \end{frame}
    %-------------------------------------------------------------------------------------------

%-----------------------------------------------------------------------------------------------
\subsection{Relações para Calores Específicos}
%-----------------------------------------------------------------------------------------------

    % !j 96 -i8
    %-------------------------------------------------------------------------------------------
    \begin{frame}[allowframebreaks]{Relações para Calores Específicos --}\vspace*{-0em}
        Tomando-se \alert{$(\partial M/\partial y)_x = (\partial N/\partial x)_y$}
        com \alert{$dz = M\,dx + N\,dy$} partindo de

        \begin{equation*}
            ds = 
                \frac{c_v}{T}\,dT +
                \left(
                    \frac{\partial P}{\partial T}
                \right)_v\,dv,
        \end{equation*}

        \noindent tem-se:

        \begin{equation*}
            \alert{
                \left(
                    \frac{\partial c_v}{\partial v}
                \right)_T =
                T\left(
                    \frac{\partial^2 P}{\partial T^2}
                \right)_v
            }.
        \end{equation*}

        \pagebreak
        Analogamente, partindo-se de

        \begin{equation*}
            ds = 
                \frac{c_p}{T}\,dT -
                \left(
                    \frac{\partial v}{\partial T}
                \right)_P\,dP,
        \end{equation*}

        \noindent chega-se a:

        \begin{equation*}
            \alert{
                \left(
                    \frac{\partial c_p}{\partial P}
                \right)_T =
                -T\left(
                    \frac{\partial^2 v}{\partial T^2}
                \right)_P
            }.
        \end{equation*}

        \pagebreak
        Eliminando-se \alert{$ds$} de \alert{$d[s\!:\!s(T, v)] = d[s\!:\!s(T, P)]$}, chega-se a:

        \begin{equation*}
            dT =
                \frac{T(\partial P/\partial T)_v}{c_p - c_v}\,dv +
                \frac{T(\partial v/\partial T)_P}{c_p - c_v}\,dP;
        \end{equation*}\vspace*\medskipamount

        \noindent tomando-se \alert{$d[T\!:\!T(v, P)]$}, e equacionando coeficientes de
        \alert{$dv$} OU \alert{$dP$}, dá:

        \begin{equation*}
            \alert{
                c_p - c_v =
                    T\left(
                        \frac{\partial v}{\partial T}
                    \right)_P
                    \left(
                        \frac{\partial P}{\partial T}
                    \right)_v =
                    -T\left(
                        \frac{\partial v}{\partial T}
                    \right)_P^2
                    \left(
                        \frac{\partial P}{\partial v}
                    \right)_T
            }.
        \end{equation*}

        \pagebreak
        Sejam a \emph{razão de calores específicos}, \alert{$\gamma$} e o \emph{coeficiente de
        expansão volumétrica}, \alert{$\beta$}:

        \begin{equation*}
            \alert{\gamma \equiv \frac{c_p}{c_v}},
            \quad\mbox{e}\quad
            \alert{\beta \equiv \frac{1}{v}\left(\frac{\partial v}{\partial T}\right)_p}.
        \end{equation*}\vspace*\medskipamount

        Sejam a \emph{compressibilidade isotérmica}, \alert{$\kappa_T$}, e \emph{isentrópica},
        \alert{$\kappa_s$}:

        \begin{equation*}
            \alert{\kappa_T \equiv \frac{-1}{v}\left(\frac{\partial v}{\partial P}\right)_T},
            \quad\mbox{e}\quad
            \alert{\kappa_s \equiv \frac{-1}{v}\left(\frac{\partial v}{\partial P}\right)_s}.
        \end{equation*}\vspace*\medskipamount

        As propriedades \alert{$(\beta, \kappa_T, \kappa_s)$} são mensuráveis à partir de
        \alert{$(P, T, v)$} e suas variações.

        \pagebreak
        Então, a diferença de calores específicos pode ser expressa pela \alert{relação de
        Mayer}:

        \begin{equation*}
            \alert{
                c_p - c_v =
                    \frac{vT\beta^2}{\kappa_T}
            },
        \end{equation*}

        \noindent pela qual deduz-se:\vspace*\medskipamount

        \begin{enumerate}
            \item \alert{$c_p \geqslant c_v$}, por $v$, $T$, $\kappa_T$ e $\beta^2$ serem
                positivos;
            \item \alert{$c_p - c_v \to 0$ à medida em que $T \to 0$};
            \item \alert{$c_p = c_v$ para substâncias incompressíveis}, $(dv = 0)$.
        \end{enumerate}

    \end{frame}
    %-------------------------------------------------------------------------------------------

%-----------------------------------------------------------------------------------------------
\section{Relações de Outras Propriedades}
%-----------------------------------------------------------------------------------------------

%-----------------------------------------------------------------------------------------------
\subsection{Relações Isentrópicas --- $k$ e $\mathfrak{c}$}
%-----------------------------------------------------------------------------------------------

    % !j 96 -i8
    %-------------------------------------------------------------------------------------------
    \begin{frame}[allowframebreaks]{Expoente de Expansão Isentrópica --}\vspace*{-0em}
        Para \alert{$Pv^k = \mbox{const.}$}~em processo isentrópico, \alert{$P$} e \alert{$v$}
        devem seguir:

        \begin{equation*}
            0 = \frac{dP}{P} + k\frac{dv}{v}.
        \end{equation*}

        Tomando-se \alert{$s\!:\!s(P, v)$}, para processo isentrópicos, \alert{$ds = 0$},
        tem-se:

        \begin{equation*}
            ds = 0 =
                \left(
                    \frac{\partial s}{\partial P}
                \right)_v\,dP +
                \left(
                    \frac{\partial s}{\partial v}
                \right)_P\,dv = \alert{
                -\left(
                    \frac{\partial v}{\partial T}
                \right)_s\,dP +
                \left(
                    \frac{\partial P}{\partial T}
                \right)_s\,dv
            },
        \end{equation*}

        \pagebreak
        Tornando o coeficiente de \alert{$dP$} unitário, a definição apropriada de \alert{$k$}
        é:

        \begin{equation*}
            \alert{
                k \equiv
                    \frac{-v}{P}
                    \left(
                        \frac{\partial P}{\partial v}
                    \right)_s
            }.
        \end{equation*}\vspace*\medskipamount

        Usando as relações \alert{cíclica}, de \alert{reciprocidade} e todas as \alert{4 de
        Maxwell}, chega-se em:

        \begin{equation*}
            \alert{
                k =
                    \frac{-v}{P}
                    \left(
                        \frac{\partial P}{\partial v}
                    \right)_T
                    \frac{c_p}{c_v}
            },
        \end{equation*}

        \noindent mostrando que \alert{em geral, $k \neq \gamma$}.

        \pagebreak
        Segue uma expressão para \alert{$k$} empregando processos a \alert{$(v, P)$-const.}:

        \begin{equation*}
            \alert{
                k =
                    \frac{v}{P}
                    \frac{
                        \left(\dfrac{\partial P}{\partial T}\right)_{v}
                    }{
                        \left(\dfrac{\partial s}{\partial T}\right)_{v}
                    }
                    \left(\frac{\partial s}{\partial v}\right)_{P} =
                    \frac{vT\beta}{P\kappa_Tc_v}
                    \left(\frac{\partial s}{\partial v}\right)_{P}
            }.
        \end{equation*}

    \end{frame}
    %-------------------------------------------------------------------------------------------

    % !j 96 -i8
    %-------------------------------------------------------------------------------------------
    \begin{frame}[allowframebreaks]{Velocidade Adiabática do Som --}\vspace*{-0em}
        Seja \alert{$\mathfrak{c}$} a vlert{elocidade adiabática do som} na substância, em
        \alert{\meter\per\second}. É possível mostrar que:

        \begin{equation*}
            \alert{
                \frac{\mathfrak{c}^2}{1000} \equiv
                \left(\frac{\partial P}{\partial\rho}\right)_s =
                \gamma\left(\frac{\partial P}{\partial\rho}\right)_T
            },
        \end{equation*}\vspace*\medskipamount

        \noindent onde \alert{$\rho \equiv \nicefrac{1}{v}$}, tal que:

        \begin{equation*}
            \alert{
                \frac{\mathfrak{c}^2}{1000} =
                -v^2 \left(\frac{\partial P}{\partial v}\right)_s =
                -\gamma v^2 \left(\frac{\partial P}{\partial v}\right)_T
            }.
        \end{equation*}\vspace*\medskipamount

        O fator de \alert{$1000$} advém da conversão \alert{$\unit{1}{\kilo\joule\per\kilogram}
        = \unit{1000}{\squaren\meter\per\squaren\second}$}.

    \end{frame}
    %-------------------------------------------------------------------------------------------

%-----------------------------------------------------------------------------------------------
\subsection{Coeficientes $\mu$ na forma $(\partial T/\partial P)_{\phi}$}
%-----------------------------------------------------------------------------------------------

    % !j 96 -i8
    %-------------------------------------------------------------------------------------------
    \begin{frame}[allowframebreaks]{Coefficiente de Joule-Thomson $(\phi \equiv h)$}\vspace*{-0em}
    \end{frame}
    %-------------------------------------------------------------------------------------------

    % !j 96 -i8
    %-------------------------------------------------------------------------------------------
    \begin{frame}[allowframebreaks]{Coefficiente de Expansão Isentrópica $(\phi \equiv s)$}\vspace*{-0em}
    \end{frame}
    %-------------------------------------------------------------------------------------------

    %------------------------------------------------------------------------------------------
    \begin{frame}[allowframebreaks]{Tópicos de Leitura --}
        \begin{thebibliography}{Çengel, Y.~A., 2013}
            %-----
            \bibitem[Çengel, Y.~A., 2013]{2013-CengelYA+BolesMA-AMGH}
                Çengel, Y.~A. e Boles, M.~A.
                \newblock{%
                    {\em Termodinâmica $7^\mathrm{a}\!$ Edição\/}.
                    \alert{Seções~12-4 a 12-5.}
                }
                \newblock{\footnotesize AMGH. Porto Alegre. ISBN 978-85-8055-200-3.}
            %-----
            \bibitem[Naaktgeboren, C., 2018]{2018-NaaktgeborenC-UTFPR}
                Naaktgeboren, C.
                \newblock{%
                    {\em Thermodynamic Properties Relations (Handout)}.
                    \alert{Seções 5 e 6.}
                }
                \newblock{\footnotesize Disponibilizado no AVA.}
        \end{thebibliography}
    \end{frame}
    %------------------------------------------------------------------------------------------

    % Finishes with stunning image, with credit
    \ArtEndWSE{pexels-eugene-dorosh-739411}{jpg}{txt}

%-----------------------------------------------------------------------------------------------
\end{document}
%-----------------------------------------------------------------------------------------------

